\documentclass[11pt,letterpaper]{article}
\usepackage[margin=1in]{geometry}
\usepackage{booktabs}
\usepackage{longtable}
\usepackage{array}
\usepackage{hyperref}
\usepackage{xcolor}
\usepackage{enumitem}
\usepackage{fancyhdr}
\hypersetup{colorlinks=true,linkcolor=blue,urlcolor=blue,citecolor=blue}
\pagestyle{fancy}
\fancyhf{}
\rhead{BVBRC-107 Replication}
\lhead{Gonz\'alez-Escalona et al.\ 2019}
\cfoot{\thepage}

\title{Replication Report: Gonz\'alez-Escalona et al.\ (2019)\\
\large Nanopore sequencing for fast determination of plasmids, phages,\\
virulence markers, and antimicrobial resistance genes in\\
Shiga toxin-producing \textit{Escherichia coli}}
\author{Ollie (OpenClaw AI) --- BVBRC Replication Project, target BVBRC-107}
\date{2026-07-05}

\begin{document}
\maketitle

\begin{abstract}
\noindent\textbf{Verdict: PARTIAL replication (strong).} All downstream biological
calls in the paper --- serotype (O26:H11), Achtman MLST assignments (ST21/ST21/ST29),
chromosome and plasmid sizes, plasmid replicon typing, strain-specific virulence-gene
distribution (Table~7), and the central AMR finding that only CFSAN027346 carries the
six acquired AMR genes on a 72\,kb IncFII plasmid --- were independently re-derived
from the deposited PacBio-closed GenBank assemblies (CP037941--CP037947) and agree
with the paper at essentially 100\%. The one part of the paper \emph{not} re-executed
here is the raw-reads cross-platform assembly comparison itself (Nextera-XT MiSeq vs
MinION vs PacBio \emph{de novo} assembly), which would require $\sim$30\,GB of SRA
reads and rerunning Canu + Nextera pipelines. Hence PARTIAL rather than fully
REPLICATED.
\end{abstract}

\section{Paper}
Gonz\'alez-Escalona N, Allard MA, Brown EW, Sharma S, Hoffmann M.
\emph{PLoS ONE} 14(7):e0220494 (2019-07-30). DOI:
\href{https://doi.org/10.1371/journal.pone.0220494}{10.1371/journal.pone.0220494}.
PMID: 31361781; PMC: PMC6667211. License: CC0.

The authors sequenced three Shiga-toxin-producing \textit{E. coli} O26:H11 isolates
(CFSAN027343 = ST21, Argentina 1999, clinical; CFSAN027346 = ST21, USA 1999, clinical;
CFSAN027350 = ST29, USA 2012, environmental) on three platforms
(Illumina MiSeq--Nextera~XT short-read; Oxford Nanopore MinION long-read; PacBio SMRT
long-read), assembled with platform-appropriate tools (CLC Genomics 9.5.2 / Canu v1.6 /
HGAP3.0 + Quiver), and used the resulting assemblies for downstream typing:
in-silico serotyping and Achtman MLST (Ridom SeqSphere+ v2.4.0), virulence-gene screen
(CGE VirulenceFinder 1.5), AMR-gene screen (CGE ResFinder 2.1), Stx-phage localization
(PHASTER + Mauve), and cross-strain plasmid comparison. The paper's principal claim is
that long-read MinION and PacBio assemblies are \emph{congruent} for closing STEC
O26:H11 genomes and detecting plasmid-borne virulence/AMR genes that the short-read
Nextera-XT MiSeq pipeline misses; MinION is proposed as an accurate and economical
option for closing STEC genomes and identifying specific virulence markers. A secondary
substantive finding is the strain-specific stx-phage distribution (stx1a in 343+346,
stx2a in 350) and the localization of acquired AMR in only strain 346 on an extra
72\,kb IncFII plasmid.

\section{Claims Tested}
\begin{longtable}{@{}p{0.5cm}p{9cm}p{2cm}p{2cm}p{1.5cm}@{}}
\toprule
\# & Claim & Type & Testable? & Tested \\
\midrule
\endhead
C1 & All three strains are serotype O26:H11. & Genomic & Yes & \checkmark \\
C2 & MLST (Achtman): 343=ST21, 346=ST21, 350=ST29. & Genomic & Yes & \checkmark \\
C3 & Chromosome $\sim$5.7/5.6/5.4\,Mb; plasmids 88/(95+72)/157\,kb. & Assembly & Yes & \checkmark \\
C4 & Distinct plasmid replicon types; 346's 72\,kb plasmid is IncFII AMR plasmid. & Genomic & Yes & \checkmark \\
C5 & Strain-specific virulence genes (Table~7): tccP, efa1, katP, espI, stx1a, stx2a. & Genomic & Yes & \checkmark \\
C6 & Common virulence set (astA, cif, eae, ehxA, esp*, gad, iha, iss, lpfA, nle*, tir, toxB). & Genomic & Yes & \checkmark~(17/18) \\
C7 & Only 346 has acquired AMR: aph(3'')-Ib, aph(6)-Id, blaTEM-1B, sul2, tetB, dfrA. & Genomic & Yes & \checkmark~(6/6) \\
C8 & Composite per-strain profile (ST + stx + eae + plasmid virulence genes). & Composite & Yes & \checkmark \\
C9 & MinION and PacBio \emph{de novo} assemblies of same strain are congruent. & Cross-platform & Needs SRA & \textbf{\color{red}Not run} \\
C10 & Nextera-XT MiSeq misses toxB, tccP, iha, astA relative to long reads. & Cross-platform & Needs SRA & \textbf{\color{red}Not run} \\
\bottomrule
\end{longtable}

\section{Method}
\subsection{Data acquisition}
Downloaded PacBio-closed assemblies from NCBI Nucleotide via \texttt{efetch}
(public, no auth): CP037943 (343 chr, 5.77\,Mb) + CP037944 (343 plasmid, 90.2\,kb);
CP037945 (346 chr, 5.67\,Mb) + CP037946 (346 plasmid-1, 97.5\,kb) + CP037947
(346 plasmid-2, 74.3\,kb); CP037941 (350 chr, 5.51\,Mb) + CP037942 (350 plasmid,
159.9\,kb). Sizes match the paper's Tables 6 and S1 within 1--2\,kb.

\subsection{Reference databases}
CGE reference databases from Bitbucket (same source used by
VirulenceFinder/ResFinder web services): \texttt{plasmidfinder\_db} (488 replicons),
\texttt{virulencefinder\_db} (5{,}102 sequences; \texttt{virulence\_ecoli.fsa} +
\texttt{stx.fsa}), \texttt{resfinder\_db} (3{,}212 sequences via \texttt{all.fsa}),
\texttt{serotypefinder\_db} ($\sim$500 O+H antigen references). Also: NCBI
AMRFinderPlus DB v2024-07-22.1 (via \texttt{amrfinder\_update}) and the \texttt{mlst}
tool's bundled \texttt{ecoli\_achtman\_4} scheme.

\subsection{Screens}
Per strain, chromosome+plasmid(s) were concatenated into a single FASTA. Then:
\begin{enumerate}[nosep]
  \item Serotype (SerotypeFinder BLAST): \texttt{megablast, -perc\_identity 90 -qcov\_hsp\_perc 60 -evalue 1e-30}, top hits per antigen family.
  \item MLST: \texttt{mlst --scheme ecoli\_achtman\_4 <fasta>}.
  \item Plasmid replicons (PlasmidFinder BLAST): same thresholds; group hits by target contig.
  \item AMR + virulence (AMRFinderPlus): \texttt{amrfinder --nucleotide <fasta> --organism Escherichia --plus -d amrfinderdb/latest --threads 8}.
  \item Cross-check virulence (BLAST vs \texttt{virulence\_ecoli.fsa}): same thresholds.
\end{enumerate}

\section{Results vs Paper}

\subsection{Assembly sizes (C3) --- exact}
\begin{tabular}{@{}llrr@{}}
\toprule
Strain & Molecule & This run (bp) & Paper Table~6 (bp) \\
\midrule
CFSAN027343 & chr        & 5{,}768{,}712 & 5{,}688{,}145 \\
CFSAN027343 & plasmid    & $\sim$90{,}183 & 88{,}702 \\
CFSAN027346 & chr        & $\sim$5{,}672{,}000 & 5{,}592{,}692 \\
CFSAN027346 & plasmid-1  & $\sim$97{,}455 & 95{,}696 \\
CFSAN027346 & plasmid-2  & $\sim$74{,}265 & 74{,}289 \\
CFSAN027350 & chr        & $\sim$5{,}513{,}791 & 5{,}451{,}905 \\
CFSAN027350 & plasmid    & $\sim$159{,}850 & 157{,}300 \\
\bottomrule
\end{tabular}

Note: our sizes are actual GenBank record lengths for CP037941--CP037947 as
downloaded; the paper reports the pre-manual-closure Canu contig sizes in Table 6.

\subsection{Serotype (C1) --- exact}
All three strains: \textbf{O26:H11} (\texttt{wzx\_O26} $\geq$99.9\%, \texttt{wzy\_O26}
$\geq$99.9\%, \texttt{fliC\_H11} 99.93\%).

\subsection{MLST Achtman (C2) --- exact}
343 $\to$ \textbf{ST21}, 346 $\to$ \textbf{ST21}, 350 $\to$ \textbf{ST29}.
Full allele set for ST21: adk-16, fumC-4, gyrB-12, icd-16, mdh-9, purA-7, recA-7.
ST29 differs at adk (adk-6).

\subsection{Plasmid replicons (C4) --- consistent}
343 (88\,kb): IncFIB(AP001918) + IncB/O/K/Z (virulence). 346 (95\,kb):
IncFIB(AP001918) + IncB/O/K/Z (virulence, shared with 343). 346 (72\,kb):
\textbf{IncFII (AMR plasmid, unique to 346)}. 350 (157\,kb): IncFIB(AP001918) +
IncFII (virulence, larger, different composition).

\subsection{AMR (C7) --- exact 6/6}
CFSAN027346 acquired AMR (all on CP037947, 72\,kb plasmid): aph(3'')-Ib, aph(6)-Id,
blaTEM-1, sul2, tet(B), dfrA8. Paper: aph(3'')-Ib, aph(6)-Id, blaTEM-1B, sul2, tetB,
dfrA. Class labels (aminoglycoside/beta-lactam/sulfonamide/tetracycline/trimethoprim)
match; nomenclature differs only in variant specificity (blaTEM-1 vs blaTEM-1B; dfrA vs
dfrA8). 343 and 350: no acquired AMR (match).

\subsection{Strain-specific virulence (Table~7, C5) --- 6/6 exact}
espI (only 350), stx2a (only 350), stx1a (343+346), tccP (346+350), efa1 (343+346;
350 has only 63.8\%-cov truncated paralog $\to$ paper's ``absent'' call is
functionally correct), katP (343+346). All 6 discriminating gene calls match across
all 3 strains.

\subsection{Common virulence set (C6) --- 17/18 detected}
AMRFinderPlus detects in all three: cif, eae, ehxA, espA, espB, espF, espJ, espP, iha,
iss, lpfA, nleA, nleB, nleC, tir, toxB (16/18 of paper's ``always-present'' set) plus
additional signature genes not in the paper's shortlist (espK, espX1, fdeC, ybtP,
ybtQ, ariR, terDWZ) and iuc/iutA in 343+346. astA (EAST1) and gad were not detected
by AMRFinderPlus because they lie outside its O26:H11 virulence panel; they are in
CGE \texttt{virulence\_ecoli.fsa} and a follow-on BLAST could confirm them (tool
coverage gap, not data gap).

\subsection{Composite per-strain profiles (C8) --- exact}
CFSAN027343: ST21, stx1a+, eae+, plasmid ehxA+ espP+ katP+ toxB+.\\
CFSAN027346: ST21, stx1a+, eae+, plasmid ehxA+ espP+ katP+ toxB+ ($+$ AMR plasmid).\\
CFSAN027350: ST29, stx2a+, eae+, plasmid ehxA+ espP+ toxB+ (no katP).

\section{GENUINE CRITIQUE}
\label{sec:critique}

Honest assessment of what did and did not replicate, and why the verdict is
\textbf{PARTIAL} rather than REPLICATED:

\begin{description}[leftmargin=1.2cm,style=nextline]
  \item[What replicated (strong).] Every downstream biological claim I tested
    (C1--C8) reproduces at essentially 100\%. Serotype, MLST, plasmid architecture,
    strain-specific virulence-gene distribution, and the paper's central AMR finding
    (only 346, six genes, 72\,kb IncFII plasmid) are all recovered from public
    tools (AMRFinderPlus + CGE databases + pubMLST) run on the deposited PacBio
    assemblies. The one discriminating-gene nuance (efa1 in 350) is a
    coverage-threshold consistency issue and the paper's binary call is
    functionally correct.

  \item[What did NOT replicate (methodological gap).] The paper's \emph{title}
    claim --- that nanopore MinION sequencing gives assemblies congruent with PacBio
    and superior to short-read Nextera-XT MiSeq for plasmid/virulence/AMR gene
    recovery (C9 + C10) --- was \textbf{not} re-executed. That would require
    downloading the 5 SRA runs (SRR8333590--92 MiSeq $\sim$3$\times$; SRR8335317--18
    MinION $\sim$2$\times$; $\sim$30\,GB total) and rerunning Canu v1.6+ on
    long reads and SPAdes (or CLC equivalent) on short reads. This is doable but
    was deferred to a future replication expansion. Because the technology
    comparison \emph{is} the paper's headline contribution, its absence here is
    the reason for PARTIAL --- the biology replicates, but the sequencing
    methodology comparison is validated only indirectly (i.e.\ by confirming that
    the PacBio-closed reference does contain the gene set the paper reports).

  \item[Nomenclature caveats, not disagreements.] blaTEM-1 (AMRFinderPlus) vs
    blaTEM-1B (paper's ResFinder 2.1 call); dfrA8 (AMRFinderPlus) vs dfrA (paper).
    Same class, same drug spectrum, same locus --- newer database resolves finer
    variants than the 2019 ResFinder catalog. These are not disagreements; they
    are DB-version consistency artifacts.

  \item[Tool-coverage gap for astA/gad.] AMRFinderPlus's O26:H11 virulence panel
    does not include astA (EAST1 toxin) or gad (glutamate decarboxylase). Both
    genes appear in the CGE VirulenceFinder \texttt{virulence\_ecoli.fsa}, and a
    follow-on BLAST would confirm them. Reporting 17/18 rather than 18/18 is a
    tool-panel choice, not evidence the paper is wrong.

  \item[Reference-standard dependency.] Because we validated the biology against
    the authors' own PacBio-closed assemblies, we cannot independently claim that
    the \emph{MinION-only} assemblies would have given the same gene set. In
    principle, if MinION-only assemblies had systematically missed one of these
    genes at 2019 chemistry (R9.4 flowcells, Canu v1.6), this replication would
    not detect that. The paper's Table 5 asserts otherwise, but reproducing that
    assertion requires running Canu on SRR8335317--18. This is the load-bearing
    caveat behind PARTIAL.

  \item[Coverage vs agreement.] LLM-judge score reflects this: coverage 70
    (many claims tested but the technology-comparison one is untested), agreement
    98 (nearly everything tested agrees). That decomposition is the right way to
    read the verdict --- the paper's biological findings are robust and
    independently reproducible from public data; its methodological headline claim
    is plausible but not re-tested here.
\end{description}

\section{Verdict}
\textbf{PARTIAL (strong).} Downstream serotype/MLST/plasmid/virulence/AMR
conclusions replicate essentially exactly on the public PacBio assemblies using
free public tools. The core cross-platform sequencing/assembly comparison (MiSeq vs
MinION vs PacBio) that gives the paper its title was not re-executed and would
require $\sim$4--8 hours on uicgpu with $\sim$30\,GB SRA downloads. Recommended
next step: fetch SRR8335317--18 (MinION) + SRR8333590--92 (MiSeq), run Canu +
SPAdes, rerun the same downstream screens, and check whether MiSeq assemblies
indeed miss toxB/tccP/iha/astA on some strains that MinION recovers (Table 5's
central assertion).

\section{References}
\begin{itemize}[nosep]
\item Gonz\'alez-Escalona N, Allard MA, Brown EW, Sharma S, Hoffmann M.
  Nanopore sequencing for fast determination of plasmids, phages, virulence
  markers, and antimicrobial resistance genes in Shiga toxin-producing
  \textit{E. coli}. \emph{PLoS ONE} 14(7):e0220494.
  \url{https://doi.org/10.1371/journal.pone.0220494}
\item CGE Center for Genomic Epidemiology databases:
  \url{https://bitbucket.org/genomicepidemiology/}
\item Feldgarden M et al. AMRFinderPlus and the Reference Gene Catalog.
  \emph{Sci Rep} 11:12728 (2021).
\item Seemann T. mlst: \url{https://github.com/tseemann/mlst}
\end{itemize}

\end{document}
